\section{Introduction}

Imagine a visitor standing in a virtual classroom, pointing at a skeleton, and
asking an embodied agent, ``What is this?'' The visible object, the visitor's
selection, the responsible agent, and the relevant knowledge are all part of the
interaction. A conversational response is useful only if the system resolves
the intended object, selects an agent that is responsible for it, applies an
appropriate capability, and makes failures recoverable. The same requirement
holds when the object is moved, renamed, or assigned to another agent.

This apparently simple interaction crosses two development worlds. Unity
contains rooms, zones, assets, colliders, avatars, and interaction components.
The agent backend contains personas, prompts, knowledge, routing rules, and
language-model calls. In a directly wired prototype, the relationships between
these elements are encoded in several files and pieces of code. A response may
be technically successful while being grounded in the wrong object or produced
by the wrong agent. When the scene changes, stale relationships can remain
hidden until a person encounters a failure at runtime.

We investigate whether an explicit model-grounding layer can make these
relationships executable and observable. Our system, \systemname{}, represents
user-facing scenarios, spatial targets, embodied agents, capabilities, runtime
bindings, expected assertions, and observed outcomes in a shared interaction
model. A generation pipeline materializes this model in a Unity/WebXR client and
a multi-agent backend. Runtime events carry model identifiers back into a
validation layer, allowing a developer to follow an interaction from a user
request to its outcome.

This problem sits at a junction that prior systems usually address in parts.
Embodied agents make language and social behavior visible
\cite{cassell1999embodiment,kopp2005museum}; situated-language work grounds
references in shared environments
\cite{gorniak2005grounding,lemaignan2012grounding}; conversational XR tools
support agent integration or end-user authoring
\cite{delatorre2024llmr,carcangiu2025tellxr,buldu2025cuify,liu2026agenthands};
and model-based interface engineering separates interaction concepts from
platform realizations
\cite{puerta1999framework,calvary2003cameleon,paterno2009maria}. Our question is
not whether any one of these ingredients is new. It is whether explicitly
connecting them provides a useful control and evidence boundary for a modern
spatial multi-agent interface.

The model is not presented as an end in itself. It is infrastructure for an
intelligent user interface: it determines what a spatial request refers to,
which embodied agent may act, which operation is invoked, and what evidence can
explain the result. This focus distinguishes the work from a general-purpose
software architecture or metamodel contribution.

We ask three questions:

\begin{description}
  \item[RQ1: Interaction correctness.] How accurately does model-grounded
  execution resolve spatial targets and route requests across heterogeneous
  scenes, including ambiguous, negative, and changed-scene cases?

  \item[RQ2: Robust authoring and evolution.] Compared with direct artifact
  wiring, how does the model-grounded workflow affect consistency, detected
  faults, and the number of artifacts that must be edited when an interaction
  changes?

  \item[RQ3: Runtime accountability.] To what extent can an interaction be
  traced from intent and spatial target through agent capability and runtime
  action to a validated outcome?
\end{description}

This paper makes four contributions:

\begin{enumerate}
  \item a model-grounded authoring architecture for spatial embodied-agent
  interfaces, implemented as a Unity desktop prototype with a WebXR deployment
  path and a conversational-agent backend;
  \item an executable interaction chain linking scenarios, targets, agents,
  capabilities, runtime actions, assertions, and evidence;
  \item a reproducible benchmark over multiple environments covering reference
  resolution, routing, authoring changes, trace completeness, and injected
  faults; and
  \item a controlled comparison with direct artifact wiring, with claims bounded
  to the measured interaction and engineering outcomes.
\end{enumerate}

Our current evaluation is a system and computational evaluation, not a claim
about preference, trust, workload, or perceived usability. We make this boundary
explicit because technical observability can support a human-centered interface
without, by itself, proving a better user experience.
